\section{Background}
\label{sec:background}

% Draft 26 Oct 2026. Defines only what Sections 3 and 4 use.

\paragraph{Path patterns in \sqlpgq.} A property graph in \sqlpgq is a view over tables. \texttt{CREATE PROPERTY
GRAPH} names vertex tables and edge tables and says which edge columns reference which vertex keys. A \texttt{MATCH}
clause inside \texttt{GRAPH\_TABLE} binds a pattern against it. A \emph{path pattern} such as
\texttt{p = ANY SHORTEST (a)-[k:Knows]->\{1,8\}(b)} asks for one shortest path of one to eight \texttt{Knows} edges
between every pair of vertices $a$ and $b$ that the rest of the pattern admits; the quantifier bounds the length, and
\texttt{ANY SHORTEST} fixes which of the possibly many paths is returned. Four functions expose the path,
\texttt{path\_length(p)}, \texttt{vertices(p)}, \texttt{edges(p)} and \texttt{element\_id(p)}; the last is the list of
vertex and edge identifiers in path order. Everything in this paper concerns these patterns on a single
edge table, directed or undirected. In \duckpgq's rewrite of \texttt{MATCH} into SQL, vertex and edge patterns become
joins over the vertex and edge tables and the pattern's predicates become filters; that part of the rewrite is
unchanged by this work. Only the variable-length and shortest-path patterns change. They become the marker function
that the optimizer rule replaces by the path-finding operator (Section~\ref{sec:planner}).

\paragraph{\duckdb, pipelines and extensions.} \duckdb executes a query plan as pipelines that push one \duckdb
vector, 2048 rows, at a time from a source through operators into a sink; a sink that must see all its input, a hash join build or a
sort, ends one pipeline and starts the next. Pipelines are broken into tasks that a fixed pool of worker threads
executes, and a pipeline's completion is an \emph{event} whose callback can schedule further events, which is how
operators run multi-phase work without blocking a thread. \duckdb has an extension mechanism through which external code
adds to the system behind defined interfaces, among them parser extensions, scalar and table functions, optimizer
extensions, custom operators and index types (Section~\ref{sec:hooks} lists them). \duckpgq is such an extension. It adds \sqlpgq
parsing, the path-finding operator and the CSR index through these interfaces, and it gets nothing else. The storage
format, the transaction manager and the executor are \duckdb's and are unchanged.

\paragraph{MVCC and row visibility.} \duckdb keeps multiple versions of a row and gives each transaction a
\emph{visibility bound}. A row inserted by a commit before the bound, and not deleted by a commit before it, is
visible. The check is \texttt{CanFetch(transaction, row)}; a table scan applies it to every row it reads, and
an index such as \duckdb's ART returns row identifiers without checking, because the scan that fetches the rows
checks for it. Versions are garbage-collected once no active transaction can observe the earlier state, so
visibility can be asked about a live transaction only. Row identifiers are stable while a table carries any index
other than an ART. \todo{Audit 28 Oct: an identifier can be handed out again after a failed commit reverts its
append, and after a checkpoint drops a fully deleted trailing row group; qualify, or show that the index handles both.}

\paragraph{Compressed sparse row.} A CSR stores a graph with $V$ vertices and $E$ edges in two arrays. The vertices are
numbered $0$ to $V-1$; an offset array of $V+1$ entries gives, for vertex $v$, the range $[\mathit{off}[v],
\mathit{off}[v+1])$ in an edge array of $E$ destination ids sorted by source. The neighbours of $v$ are one contiguous
slice, a traversal reads the edge array sequentially, and the whole structure costs $4$ to $8$ bytes per edge. Its
weakness is that an insert in the middle moves everything behind it, so a CSR is built once from a sorted edge list
and rebuilt when the graph changes; Section~\ref{sec:index} is about keeping one under updates without rebuilding it
per query. A \emph{sparse-row} CSR keeps offsets only for the sources that have edges, as a sorted list of (source,
offset) pairs, which matters when the sources are a small subset of a large id space.

\paragraph{Multi-source BFS.} A breadth-first search from one source visits each vertex at most once and each edge
once; the cost of a level is the frontier plus its edges, but the memory access is one random read per edge.
Multi-source BFS~\cite{then2014msbfs} runs $w$ searches together, with $w$ the width of a machine word or a SIMD
register, and keeps for each vertex a bit set of the searches that have seen it and a bit set of the searches whose
frontier it is on. Relaxing an edge $u \to v$ becomes $\mathit{next}[v] \mathrel{|}= \mathit{visit}[u]$, one read and one write
of $w$ bits for all $w$ searches at once, and a search terminates when its destination's bit is set. Its cost per
level is a sweep over all vertices whose visit bits are non-zero, which is cheap per search when many searches are
live and expensive when one is. \duckpgq runs 512 lanes per wave, the width of eight 64-bit words.
